\begin{enumerate}
\item $x = 600 - 450 = 150$

$y = 750 - 450 = 300$

Dans le contexte de l'exercice :
\begin{itemize}
	\item $x$ : il y a 150 graines de basilic qui germent.
	\item $y$ : il y a 300 graines d'estragon qui ne germent pas.
\end{itemize}


\item
On considère que chaque graine a la même probabilité d'être semée.

\textbf{Hypothèse d'équiprobabilité} : les proportions sont assimilables à des probabilités.

\begin{enumerate}
	\item Il y a \np{1000} graines, dont 600 de basilic, donc :\quad
	$P(B) = \dfrac{600}{\np{1000}} = \np{0.6}$.

	Et il y a 400 graines d'estragon, donc :\quad
	$P(E) = \dfrac{400}{\np{1000}} = \np{0.4}$.

	\item
$E \cap G$ est l'évènement «~La graine semée est une graine d'estragon et elle germe.~»

	100 graines sont des graines d'estragon qui germent, donc :\quad
	$P(E \cap G) = \dfrac{100}{\np{1000}} = \np{0.1}$.



	\item Parmi les 250 graines qui germent, 100 sont des graines d'estragon, donc :

	$P_{G}(E) = \dfrac{100}{250} = \dfrac{400}{\np{1000}} = \np{0.4}$.

	\item Parmi les 400 graines d'estragon, 100 germent, donc : \quad$P_{E}(G) = \dfrac{100}{400} = \dfrac{1}{4}$.

	\item
	L'affirmation du jardinier revient à dire que les évènements $G$ et $E$ sont indépendants.

	Or $P_{G}(E) = P(E)$, ce qui prouve que les évènements $G$ et $E$ \textbf{sont} indépendants (ou, de façon équivalente que $G$ et $B$ sont indépendants, car $B$ est l'évènement contraire de $E$).
	\end{enumerate}
\end{enumerate}
